\chapter{Limites, Risques et Aspects Éthiques}

\section{Limites techniques}
\begin{itemize}[leftmargin=2em]
  \item dépendance au seuil de similarité pour la décision finale,
  \item sensibilité du comportement selon le modèle choisi,
  \item coût et latence liés aux appels LLM externes.
\end{itemize}

\section{Biais et qualité des données}
Les biographies publiques peuvent refléter des biais éditoriaux, des asymétries de couverture médiatique et des formulations culturellement situées. Le pipeline hérite partiellement de ces biais même après normalisation.

\section{RGPD et données publiques}
Le projet exploite des sources publiques (Wikipédia). Même dans ce cadre, il est recommandé de:
\begin{itemize}[leftmargin=2em]
  \item éviter des usages de profilage sensible,
  \item tracer les sources,
  \item conserver une boucle de validation humaine pour les usages à enjeu.
\end{itemize}

\section{Lecture honnête des métriques}
Les performances proches de 1.00 sur certains splits ne suffisent pas à conclure à une robustesse générale. Le split hard et les baselines sont indispensables pour une interprétation crédible.

